% SPDX-License-Identifier: CC-BY-NC-ND-4.0
% !TEX root = main.tex

\section{極限}

\begin{definition} \label{def: Cone}
$\mathscr{J}, \mathscr{C}$を圏とし$F \colon \mathscr{J} \to \mathscr{C}$を関手とする.\ $F$を\textbf{形}(shape)$\mathscr{J}$の\textbf{図式}(diagram)と呼ぶ.\ $A \in \Ob(\mathscr{C})$とし$C_A \colon \mathscr{J} \to \mathscr{C}$を定関手とする.\ $A$と自然変換$\eta \colon C_A \Rightarrow F$の組$(A, \eta)$を$A$を\textbf{頂点}(apex)とする$F$への\textbf{錐}(cone)という.
% https://q.uiver.app/#q=WzAsOCxbMCwwLCJDX0EoXFxtYXRocm17ZG9tfShmKSkiXSxbMSwwLCJGKFxcbWF0aHJte2RvbShmKX0pIl0sWzEsMSwiRihcXG1hdGhybXtjb2R9KGYpKSJdLFswLDEsIkNfQShcXG1hdGhybXtjb2R9KGYpKSJdLFszLDAsIkEiXSxbMywxLCJBIl0sWzQsMCwiRihcXG1hdGhybXtkb20oZil9KSJdLFs0LDEsIkYoXFxtYXRocm17Y29kfShmKSkiXSxbMCwxLCJcXGV0YV97XFxtYXRocm17ZG9tfShmKX0iXSxbMywyLCJcXGV0YV97XFxtYXRocm17Y29kfShmKX0iXSxbMSwyLCJGKGYpIl0sWzAsMywiQ19BKGYpIiwyXSxbNCw1LCIiLDAseyJsZXZlbCI6Miwic3R5bGUiOnsiaGVhZCI6eyJuYW1lIjoibm9uZSJ9fX1dLFs2LDcsIkYoZikiXSxbNCw2LCJcXGV0YV97XFxtYXRocm17ZG9tfShmKX0iXSxbNSw3LCJcXGV0YV97XFxtYXRocm17Y29kfShmKX0iXSxbMTAsMTIsIiIsMCx7InNob3J0ZW4iOnsic291cmNlIjo0MCwidGFyZ2V0IjozMH0sInN0eWxlIjp7InRhaWwiOnsibmFtZSI6ImFycm93aGVhZCJ9fX1dXQ==
\[\begin{tikzcd}
	{C_A(\dom(f))} & {F(\dom(f))} && A & {F(\dom(f))} \\
	{C_A(\cod(f))} & {F(\cod(f))} && A & {F(\cod(f))}
	\arrow["{\eta_{\dom(f)}}", from=1-1, to=1-2]
	\arrow["{C_A(f)}"', from=1-1, to=2-1]
	\arrow[""{name=0, anchor=center, inner sep=0}, "{F(f)}", from=1-2, to=2-2]
	\arrow["{\eta_{\dom(f)}}", from=1-4, to=1-5]
	\arrow[""{name=1, anchor=center, inner sep=0}, equals, nfold, from=1-4, to=2-4]
	\arrow["{F(f)}", from=1-5, to=2-5]
	\arrow["{\eta_{\cod(f)}}", from=2-1, to=2-2]
	\arrow["{\eta_{\cod(f)}}", from=2-4, to=2-5]
	\arrow[between={0.4}{0.7}, Rightarrow, nfold, 2tail reversed, from=0, to=1]
\end{tikzcd}\]
\end{definition}

\begin{remark} \label{rem: functor_category_is_locally_small}
$\mathscr{J}$を小圏,\ $\mathscr{C}$を局所小圏とすると関手圏$\mathscr{C}^{\mathscr{J}}$は局所小圏である.
\end{remark}

\begin{definition} \label{def: coneFunctor}
$\mathscr{J}$を小圏,\ $\mathscr{C}$を局所小圏とし$F \colon \mathscr{J} \to \mathscr{C}$を形$\mathscr{J}$の図式とする.\ 関手$\Cone(-, F) \colon \mathscr{C}^{\op} \to \mathbf{Set}$を次のように定める.
\begin{itemize}
    \item $A \in \Ob(\mathscr{C})$に対して$\Cone(A, F)$を$A$を頂点とする$F$への錐全体の集合とする.
    \item $f \in \Mor(\mathscr{C})$に対して$(\cod(f), \eta) \in \Cone(\cod(f), F)$とし$\Delta \colon \mathscr{C} \to \mathscr{C}^{\mathscr{J}}$を定図式関手とする.\ このとき$\Cone(f, F)(\cod(f), \eta) \coloneq (\dom(f), \eta \circ \Delta(f))$と定める.\ つまり$g \in \mathscr{J}$に対して次の図式が可換になる.
\end{itemize}
% https://q.uiver.app/#q=WzAsNCxbMSwxLCJcXG1hdGhybXtjb2R9KGYpIl0sWzAsMiwiRihcXG1hdGhybXtkb219KGcpKSJdLFsyLDIsIkYoXFxtYXRocm17Y29kfShnKSkiXSxbMSwwLCJcXG1hdGhybXtkb219KGYpIl0sWzAsMSwiXFxldGFfe1xcbWF0aHJte2RvbX0oZyl9Il0sWzAsMiwiXFxldGFfe1xcbWF0aHJte2NvZH0oZyl9IiwyXSxbMywwLCJmIl0sWzMsMSwiKFxcZXRhIFxcY2lyYyBcXERlbHRhKGYpKV97XFxtYXRocm17ZG9tfShnKX0iLDIseyJjdXJ2ZSI6MX1dLFszLDIsIihcXGV0YSBcXGNpcmMgXFxEZWx0YShmKSlfe1xcbWF0aHJte2NvZH0oZyl9IiwwLHsiY3VydmUiOi0xfV0sWzEsMiwiRihnKSIsMl1d
\[\begin{tikzcd}
	& {\mathrm{dom}(f)} & \\
	& {\mathrm{cod}(f)} \\
	{F(\mathrm{dom}(g))} && {F(\mathrm{cod}(g))}
	\arrow["f", from=1-2, to=2-2]
	\arrow["{(\eta \circ \Delta(f))_{\mathrm{dom}(g)}}"', curve={height=6pt}, from=1-2, to=3-1]
	\arrow["{(\eta \circ \Delta(f))_{\mathrm{cod}(g)}}", curve={height=-6pt}, from=1-2, to=3-3]
	\arrow["{\eta_{\mathrm{dom}(g)}}", from=2-2, to=3-1]
	\arrow["{\eta_{\mathrm{cod}(g)}}"', from=2-2, to=3-3]
	\arrow["{F(g)}"', from=3-1, to=3-3]
\end{tikzcd}\]
この関手を\textbf{錐関手}(cone functor)という.
\end{definition}

% 任意の錐から極限となる錐に対して一意な射が存在するため錐の圏の終対象として極限が定まる
\begin{definition} \label{def: limit}
$\mathscr{J}$を小圏,\ $\mathscr{C}$を局所小圏とし$F \colon \mathscr{J} \to \mathscr{C}$を形$\mathscr{J}$の図式とする.\ $\displaystyle \int_{\mathscr{C}^{\op}} \Cone(-, F)$の終対象$(L_F, (L_F, \pi_F))$を$F$の\textbf{極限}(limit)という.
\end{definition}

\begin{definition} \label{def: IsJComplete}
$\mathscr{J}$を小圏,\ $\mathscr{C}$を局所小圏とし$F \colon \mathscr{J} \to \mathscr{C}$を形$\mathscr{J}$の図式とする.\ $F \colon \mathscr{J} \to \mathscr{C}$が極限を持つとき$\mathscr{C}$は$\mathscr{J}$\textbf{完備}($\mathscr{J}$-complete)であるという.
\end{definition}

\begin{definition} \label{def: limitFunctor}
$\mathscr{J}$を小圏とし$\mathscr{C}$を$\mathscr{J}$完備な局所小圏とする.\ 共変関手$\varprojlim_{\mathscr{J}} \colon \mathscr{C}^{\mathscr{J}} \to \mathscr{C}$を次のようにして定める.
\begin{itemize}
	\item $F \in \Ob(\mathscr{C}^{\mathscr{J}})$に対して$(L_F, (L_F, \pi_F))$を$F$の極限としたとき$\varprojlim_{\mathscr{J}}(F) \coloneq L_F$とする.
	\item $F, G \in \Ob(\mathscr{C}^{\mathscr{J}})$とし$(L_G, (L_G, \pi_G))$を$G$の極限としたとき$\eta \in \Mor_{\mathscr{C}^{\mathscr{J}}}(F, G)$に対して$(L_F, \eta \circ \pi_F) \in \Cone(L_F, G)$なのでただ1つの射$f \colon (L_F, (L_F, \eta \circ \pi_F)) \to (L_G, (L_G, \pi_G))$が存在する.\ そこで$\varprojlim_{\mathscr{J}}(\eta) \coloneq f$として定める.
\end{itemize}

$\varprojlim_{\mathscr{J}}$は共変関手になる.\ 実際$F,G,H \in \Ob(\mathscr{C}^{\mathscr{J}})$と$\eta \colon F \Rightarrow G,\ \theta \colon G \Rightarrow H$に対して$(L_F, (L_F, \eta \circ \pi_F)) \in {\displaystyle \int_{\mathscr{C}^{\op}}\Cone(-,F)}$かつ$(L_G, (L_G, \theta \circ \pi_G)) \in {\displaystyle \int_{\mathscr{C}^{\op}}\Cone(-,G)}$であり$\varprojlim_{\mathscr{J}}(\eta) \colon L_F \to L_G$と$\varprojlim_{\mathscr{J}}(\theta) \colon L_G \to L_H$は次の等式を満たす.
\vspace{0.5em}
\begin{center}
$\Cone(\varprojlim_{\mathscr{J}}(\eta), F)(L_G, \pi_G) = (L_F, \eta \circ \pi_F)$

$\Cone(\varprojlim_{\mathscr{J}}(\theta), G)(L_H, \pi_H) = (L_G, \theta \circ \pi_G)$
\end{center}
\vspace{0.5em}
また錐関手の定義(\cref{def: coneFunctor})から次の等式が成り立つ.
\vspace{0.5em}
\begin{center}
$\Cone(\varprojlim_{\mathscr{J}}(\eta), F)(L_G, \pi_G) = (L_F, \pi_G \circ \Delta(\varprojlim_{\mathscr{J}}(\eta)))$

$\Cone(\varprojlim_{\mathscr{J}}(\theta), G)(L_H, \pi_H) = (L_G, \pi_G \circ \Delta(\varprojlim_{\mathscr{J}}(\theta)))$
\end{center}
\vspace{0.5em}
したがって次の図式が可換になる.

% https://q.uiver.app/#q=WzAsNixbMCwwLCJDX3tMX0Z9Il0sWzAsMSwiQ197TF9HfSJdLFsxLDAsIkYiXSxbMSwxLCJHIl0sWzAsMiwiQ197TF9IfSJdLFsxLDIsIkgiXSxbMCwyLCJcXHBpX0YiLDAseyJsZXZlbCI6Mn1dLFsyLDMsIlxcZXRhIiwwLHsibGV2ZWwiOjJ9XSxbMSwzLCJcXHBpX0ciLDAseyJsZXZlbCI6Mn1dLFswLDEsIlxcRGVsdGEoXFx2YXJwcm9qbGltX3tcXG1hdGhzY3J7Sn19KFxcZXRhKSkiLDIseyJsZXZlbCI6Mn1dLFs0LDUsIlxccGlfSCIsMCx7ImxldmVsIjoyfV0sWzEsNCwiXFxEZWx0YShcXHZhcnByb2psaW1fe1xcbWF0aHNjcntKfX0oXFx0aGV0YSkpIiwyLHsibGV2ZWwiOjJ9XSxbMyw1LCJcXHRoZXRhIiwwLHsibGV2ZWwiOjJ9XV0=
\[\begin{tikzcd}
	{C_{L_F}} & F \\
	{C_{L_G}} & G \\
	{C_{L_H}} & H
	\arrow["{\pi_F}", Rightarrow, nfold, from=1-1, to=1-2]
	\arrow["{\Delta(\varprojlim_{\mathscr{J}}(\eta))}"', Rightarrow, nfold, from=1-1, to=2-1]
	\arrow["\eta", Rightarrow, nfold, from=1-2, to=2-2]
	\arrow["{\pi_G}", Rightarrow, nfold, from=2-1, to=2-2]
	\arrow["{\Delta(\varprojlim_{\mathscr{J}}(\theta))}"', Rightarrow, nfold, from=2-1, to=3-1]
	\arrow["\theta", Rightarrow, nfold, from=2-2, to=3-2]
	\arrow["{\pi_H}", Rightarrow, nfold, from=3-1, to=3-2]
\end{tikzcd}\]

$\Delta \colon \mathscr{C} \to \mathscr{C}^{\mathscr{J}}$は共変関手なので$\pi_H \circ \Delta(\varprojlim_{\mathscr{J}}(\theta) \circ \varprojlim_{\mathscr{J}}(\eta)) = \pi_H \circ \Delta(\varprojlim_{\mathscr{J}}(\theta)) \circ \Delta(\varprojlim_{\mathscr{J}}(\eta)) = \theta \circ \pi_G \varprojlim_{\mathscr{J}}(\eta) = (\theta \circ \eta) \circ \pi_F$となる.\ $\varprojlim_{\mathscr{J}}(\theta \circ \eta)$は次の図式を満たすただ1つの射$(L_F, (L_F, (\theta \circ \eta) \circ\pi_F)) \to (L_H, (L_H, \pi_H))$である.

% https://q.uiver.app/#q=WzAsNCxbMCwwLCJDX3tMX0Z9Il0sWzEsMCwiRiJdLFswLDEsIkNfe0xfSH0iXSxbMSwxLCJIIl0sWzEsMywiXFx0aGV0YSBcXGNpcmMgXFxldGEiLDAseyJsZXZlbCI6Mn1dLFsyLDMsIlxccGlfSCIsMCx7ImxldmVsIjoyfV0sWzAsMiwiXFxEZWx0YShcXHZhcnByb2psaW1fe1xcbWF0aHNjcntKfX0oXFx0aGV0YSBcXGNpcmMgXFxldGEpKSIsMix7ImxldmVsIjoyfV0sWzAsMSwiXFxwaV9GIiwwLHsibGV2ZWwiOjJ9XV0=
\[\begin{tikzcd}
	{C_{L_F}} & F \\
	{C_{L_H}} & H
	\arrow["{\pi_F}", Rightarrow, nfold, from=1-1, to=1-2]
	\arrow["{\Delta(\varprojlim_{\mathscr{J}}(\theta \circ \eta))}"', Rightarrow, nfold, from=1-1, to=2-1]
	\arrow["{\theta \circ \eta}", Rightarrow, nfold, from=1-2, to=2-2]
	\arrow["{\pi_H}", Rightarrow, nfold, from=2-1, to=2-2]
\end{tikzcd}\]

したがって$\varprojlim_{\mathscr{J}}(\theta) \circ \varprojlim_{\mathscr{J}}(\eta) = \varprojlim_{\mathscr{J}}(\theta \circ \eta)$である.

$\varprojlim_{\mathscr{J}}(\id_{F}) \colon L_F \to L_F$は次の等式を満たすただ1つの射である.
\begin{center}
	$\pi_F \circ \Delta(\varprojlim_{\mathscr{J}}(\id_{F})) = \id_F \circ \pi_F = \pi_F$
\end{center}
また$\Delta \colon \mathscr{C} \to \mathscr{C}^{\mathscr{J}}$は共変関手なので次の等式も成り立つ.
\begin{center}
	$\pi_F \circ \Delta(\id_{L_F}) = \pi_F \circ \id_{\Delta(L_F)} = \pi_F$
\end{center}
したがって一意性により$\varprojlim_{\mathscr{J}}(\id_{F}) = \id_{L_F} = \id_{\varprojlim_{\mathscr{J}}(F)}$となる.
\end{definition}